\section{Metode Frobenius}
\label{sec:metode_frobenius}

\noindent Untuk mendapatkan nilai frekuensi kuasinormal dari persamaan master
\eqref{eq:master-skalar}, persamaan tersebut harus diselesaikan dengan memenuhi
kondisi batas fisik pada kedua titik singularnya, yaitu di horizon peristiwa dan
di daerah terjauh. Metode Frobenius merupakan salah satu teknik penyelesaian yang 
dapat digunakan untuk menyelesaikan
persamaan diferensial linear di sekitar titik singular dengan menyatakan
solusi sebagai deret pangkat \parencite{Leaver1985, Konoplya2011, Riley2006}.

Suatu persamaan diferensial linear orde dua secara umum dapat dituliskan dalam bentuk baku
\begin{equation}
    \frac{d^2y}{dr^2} + P(r)\frac{dy}{dr} + Q(r)\,y = 0,
    \label{eq:ode-baku}
\end{equation}
dengan $P(r)$ dan $Q(r)$ merupakan koefisien yang bergantung pada $r$. Apabila kedua
koefisien tersebut bernilai berhingga dan analitik di suatu titik, titik tersebut disebut
titik biasa (\textit{ordinary point}), dan solusinya dapat dinyatakan langsung sebagai
deret Taylor biasa di sekitar titik itu \parencite{Riley2006, boas2006}.

Sebelum solusi ditentukan melalui metode Frobenius, jenis titik tempat solusi ditinjau
harus diperiksa terlebih dahulu, apakah titik tersebut merupakan titik biasa, titik singular
reguler, atau titik singular tak reguler. Pemeriksaan ini penting karena bentuk deret solusi
yang berlaku bergantung pada jenis titiknya. Pada titik biasa, solusi cukup dinyatakan
sebagai deret Taylor biasa. Pada titik singular, deret Taylor biasa tidak lagi berlaku,
sehingga metode Frobenius memperluas bentuk deret tersebut dengan mengalikannya dengan suatu
faktor pangkat $(r - r_0)^{\sigma}$,
\begin{equation}
    y(r) = (r - r_0)^{\sigma} \sum_{n=0}^{\infty} a_n (r - r_0)^n , \qquad a_0 \neq 0,
    \label{eq:frobenius-umum}
\end{equation}
dengan $r_0$ merupakan titik singular yang ditinjau dan eksponen $\sigma$ menampung perilaku
singular solusi di titik tersebut. Bentuk deret inilah yang memungkinkan metode Frobenius
menyelesaikan persamaan diferensial orde dua meskipun koefisiennya divergen, dengan syarat 
singularitasnya berjenis reguler.

Jenis suatu titik ditentukan dengan meninjau perilaku koefisien $P(r)$ dan $Q(r)$ pada
bentuk baku \eqref{eq:ode-baku} di titik tersebut. Apabila keduanya berhingga dan analitik
di suatu titik $r = r_0$, titik itu merupakan titik biasa. Sebaliknya, apabila salah satu di
antaranya divergen, titik itu merupakan titik singular yang masih dibedakan lebih lanjut
menurut tingkat kedivergenannya. Dalam menganalisis suatu titik singular $r = r_0$, ditinjau kedua limit
\begin{equation}
    p_0 = \lim_{r \to r_0} (r - r_0)\,P(r), \qquad
    q_0 = \lim_{r \to r_0} (r - r_0)^2\, Q(r).
    \label{eq:limit-singular}
\end{equation}
Apabila kedua limit pada Persamaan~\eqref{eq:limit-singular} bernilai berhingga, titik
$r = r_0$ tergolong titik singular reguler. Sebaliknya, apabila salah satu limit tersebut
divergen, titik $r = r_0$ tergolong titik singular tak reguler \parencite{Riley2006, boas2006}.

Nilai eksponen $\sigma$ pada deret Frobenius \eqref{eq:frobenius-umum} diperoleh dengan
mensubstitusikan deret tersebut ke dalam persamaan bentuk baku \eqref{eq:ode-baku}. Turunan
pertama dan kedua dari deret Frobenius adalah
\begin{equation}
    \frac{dy}{dr} = \sum_{n=0}^{\infty}(n+\sigma)\,a_n\,(r-r_0)^{n+\sigma-1}, \qquad
    \frac{d^2y}{dr^2} = \sum_{n=0}^{\infty}(n+\sigma)(n+\sigma-1)\,a_n\,(r-r_0)^{n+\sigma-2}.
    \label{eq:turunan-frobenius}
\end{equation}

Untuk mendapatkan nilai $\sigma$ yang menentukan perilaku solusi di sekitar titik singular,
ditinjau koefisien suku $n = 0$, yang khususnya hanya melibatkan $a_0$ dan $\sigma$ tanpa
bergantung pada koefisien $a_n$ lainnya. Persamaan yang diperoleh dari suku ini dikenal
sebagai persamaan indicial \parencite{Riley2006, boas2006}, yang berbentuk
\begin{equation}
    \Big[\sigma(\sigma-1) + p_0\,\sigma + q_0\Big]\,a_0 = 0.
    \label{eq:indicial-a0}
\end{equation}
Koefisien $P(r)$ dan $Q(r)$ pada bentuk baku \eqref{eq:ode-baku} tidak dapat dievaluasi
langsung di $r = r_0$ karena keduanya divergen di titik tersebut. Akan tetapi, kombinasi
$(r-r_0)P(r)$ dan $(r-r_0)^2 Q(r)$ yang muncul secara alami dari substitusi deret justru
tetap berhingga di $r = r_0$, dan nilai limitnya di titik tersebut persis $p_0$ dan $q_0$
yang telah didefinisikan pada Persamaan~\eqref{eq:limit-singular}.

Karena $a_0 \neq 0$ menurut definisi deret Frobenius, faktor di dalam kurung siku yang harus
bernilai nol, sehingga diperoleh
\begin{equation}
    \sigma(\sigma - 1) + p_0\,\sigma + q_0 = 0,
    \label{eq:indicial-umum}
\end{equation}
yang dikenal sebagai persamaan indicial \parencite{Riley2006, boas2006}. Persamaan ini
merupakan persamaan kuadrat dalam $\sigma$ sehingga memiliki dua akar, $\sigma_1$ dan
$\sigma_2$, yang menjelaskan dua perilaku solusi di titik singular tersebut.

Setelah nilai eksponen $\sigma$ diperoleh melalui persamaan indicial
\eqref{eq:indicial-umum}, ansatz Frobenius secara khusus perturbasi dapat dinyatakan secara umum sebagai
\begin{equation}
    \Psi(r) = W(r) \sum_{n=0}^{\infty} a_n\, z^n,
    \label{eq:ansatz_frobenius}
\end{equation}
dengan $z$ memetakan kondisi batas ruang-waktu menjadi $0 \le z < 1$
\parencite{Konoplya2011}. Faktor $W(r)$ menampung seluruh perilaku singular di titik-titik
batas, dan tidak selalu berupa faktor pangkat tunggal $(r-r_h)^\sigma$ seperti pada
Persamaan~\eqref{eq:frobenius-umum}. Bentuk eksplisit
$W(r)$ disusun secara terpisah pada masing-masing latar belakang.

\subsection{Metode Frobenius pada Schwarzschild}
\label{subsec:frobenius-schwarzschild}

\noindent Kerangka umum di atas diterapkan pada persamaan radial perturbasi medan skalar
\eqref{eq:pers-radial} untuk latar belakang Schwarzschild, dengan fungsi metrik
$f(r) = 1 - 2M/r$ dan horizon peristiwa terletak pada $r_h = 2M$. Membagi seluruh persamaan
dengan $f^2(r)$ membawa persamaan radial ke dalam bentuk baku \eqref{eq:ode-baku} dengan
koefisien
\begin{equation}
    P(r) = \frac{f'(r)}{f(r)}, \qquad
    Q(r) = \frac{\omega^2}{f^2(r)} - \frac{1}{f(r)}\left(\frac{\ell(\ell+1)}{r^2} + \frac{f'(r)}{r}\right),
    \label{eq:PQ-schwarzschild}
\end{equation}
dengan turunan fungsi metrik di horizon peristiwa
\begin{equation}
    f'(r_h) = \frac{2M}{r_h^2} = \frac{1}{2M}.
\end{equation}
Dengan menggunakan koefisien $P(r)$ dan $Q(r)$ pada Persamaan~\eqref{eq:PQ-schwarzschild}
serta pendekatan linear $f(r) \approx f'(r_h)(r-r_h)$ di sekitar horizon, limit pada
Persamaan~\eqref{eq:limit-singular} diperoleh sebagai
\begin{equation}
    p_0 = 1, \qquad q_0 = \frac{\omega^2}{f'(r_h)^2} = 4M^2\omega^2,
\end{equation}
sehingga persamaan indicial \eqref{eq:indicial-umum} untuk latar belakang Schwarzschild
berbentuk
\begin{equation}
    \sigma^2 + 4M^2\omega^2 = 0,
    \label{eq:indicial_horizon_saf}
\end{equation}
maka didapatkan akar-akarnya
\begin{equation}
    \sigma_{\pm} = \pm\frac{i\omega}{f'(r_h)} = \pm 2iM\omega.
    \label{eq:eksponen_horizon_saf}
\end{equation}
Kedua akar ini berpadanan dengan dua perilaku solusi di horizon peristiwa. Untuk menentukan
akar mana yang sesuai, ditinjau kembali syarat batas pada
Persamaan~\eqref{eq:bc-schwarzschild}, maka kondisi batas gelombang masuk murni memilih
$\sigma_- = -2iM\omega$.

Pada kondisi batas tak hingga spasial, sifat titik singular di $r \to \infty$ perlu ditinjau
melalui substitusi $x = 1/r$, yang memetakan $r \to \infty$ menjadi $x \to 0$, sehingga
kriteria limit \eqref{eq:limit-singular} dapat diterapkan pada titik ini. Dengan
mensubstitusikan $f(1/x) = 1 - 2Mx$ ke dalam Persamaan~\eqref{eq:PQ-schwarzschild} yang telah
ditransformasi ke variabel $x$, diperoleh
\begin{equation}
\begin{aligned}
    \tilde{P}(x) &= \frac{2}{x} - \frac{2M}{x(1-2Mx)},
    &\tilde{p}_0 &= \lim_{x\to0} x\tilde{P}(x) = 2, \\[4pt]
    \tilde{Q}(x) &= \frac{\omega^2}{x^2(1-2Mx)^2} - \frac{\ell(\ell+1)+2Mx}{x^2(1-2Mx)},
    &\tilde{q}_0 &= \lim_{x\to0} x^2\tilde{Q}(x) \to \infty.
\end{aligned}
\label{eq:transformasi-tak-hingga}
\end{equation}
Karena $\tilde{q}_0$ divergen, $x=0$ (yaitu $r\to\infty$) merupakan titik singular tak
reguler. Akibatnya, ansatz Frobenius sederhana pada
Persamaan~\eqref{eq:frobenius-umum} tidak lagi cukup, dan diperlukan \textit{generalized
Frobenius} yang mengizinkan faktor asimtotik berbentuk gabungan suku eksponensial dan suku
pangkat.

Faktor asimtotik pada tak hingga spasial diperoleh melalui koordinat \textit{tortoise}. Pada
daerah $r$ yang sangat besar, koordinat \textit{tortoise} pada
Persamaan~\eqref{eq:koordinat-tortoise} dapat diekspansi sebagai
\begin{equation}
    r_* \approx r + 2M \ln r + \text{konstan}, \qquad r \to \infty,
    \label{eq:ekspansi_tortoise_saf}
\end{equation}
sehingga solusi gelombang keluar murni $\Psi \sim e^{i\omega r_*}$ menjadi
\begin{equation}
    \Psi \sim e^{i\omega r}\, r^{2iM\omega}, \qquad r \to \infty,
    \label{eq:asimtotik_saf}
\end{equation}
yang merupakan bentuk \textit{generalized Frobenius} berupa gabungan faktor eksponensial
$e^{i\omega r}$ dan faktor pangkat $r^{2iM\omega}$, sekaligus merepresentasikan gelombang
keluar murni sesuai kondisi batas yang telah dijelaskan pada Persamaan~\eqref{eq:bc-schwarzschild}.

Faktor horizon $(r-r_h)^{-2iM\omega}$ turut menyumbang pangkat tambahan di $r \to \infty$,
karena
\begin{equation}
    (r-r_h)^{-2iM\omega} = r^{-2iM\omega}\left(1-\frac{r_h}{r}\right)^{-2iM\omega}.
    \label{eq:ekspansi-horizon-tak-hingga}
\end{equation}
Suku $(1-r_h/r)^{-2iM\omega}$ analitik dan menuju satu ketika $r\to\infty$, sehingga dapat
diserap ke dalam deret pangkat, sedangkan $r^{-2iM\omega}$ tersisa sebagai sumbangan murni di
tak hingga. Agar perilaku gelombang keluar pada Persamaan~\eqref{eq:asimtotik_saf} tetap
terpenuhi, sumbangan ini dikompensasi faktor $r^{4iM\omega}$, menghasilkan ansatz Frobenius
Schwarzschild
\begin{equation}
    \Psi(r) = e^{i\omega r}\,(r - r_h)^{-2iM\omega}\, r^{\,4iM\omega} \sum_{n=0}^{\infty} a_n\, z^n,
    \qquad z = \frac{r - r_h}{r},
    \label{eq:ansatz_saf}
\end{equation}
dengan deret $\sum a_n z^n$ reguler di seluruh domain $r_h \le r < \infty$ ($r_h = 2M$),
karena seluruh perilaku singular telah ditampung oleh faktor di depan deret.

\subsection{Metode Frobenius pada Schwarzschild Anti-de Sitter}
\label{subsec:frobenius-sads}

\noindent Kerangka umum di atas diterapkan pada persamaan radial perturbasi medan skalar
\eqref{eq:pers-radial} untuk latar belakang Schwarzschild Anti-de Sitter, dengan fungsi
metrik $f(r) = 1 - 2M/r + r^2/L^2$. Membagi seluruh persamaan dengan $f^2(r)$ membawa
persamaan radial ke dalam bentuk baku \eqref{eq:ode-baku} dengan koefisien $P(r)$ dan $Q(r)$
yang berbentuk sama seperti Persamaan~\eqref{eq:PQ-schwarzschild}, hanya dengan fungsi
metrik $f(r)$ yang berbeda.

Berbeda dengan ruang-waktu Schwarzschild, geometri asimtotik Schwarzschild Anti-de Sitter
(SAdS) memiliki batas konformal yang bersifat \textit{timelike}. Dengan sifat yang berbeda
tersebut, maka diperlukan tinjauan pada kedua batas, yaitu di sekitar horizon peristiwa
maupun di sekitar batas konformal.

Di sekitar horizon peristiwa ($r\rightarrow r_h$), pendekatan linear
$f(r) \approx f'(r_h)(r-r_h)$ yang sama seperti pada kasus Schwarzschild tetap berlaku,
sehingga limit pada Persamaan~\eqref{eq:limit-singular} kembali menghasilkan
\begin{equation}
    p_0 = 1, \qquad q_0 = \frac{\omega^2}{f'(r_h)^2}.
\end{equation}
Kedua nilai limit ini hanya bergantung pada $f'(r_h)$, bukan pada
bentuk eksplisit suku-suku $f(r)$ lainnya, sehingga persamaan indikal yang diperoleh
memiliki bentuk yang sama seperti kasus Schwarzschild,

\begin{equation}
    \sigma^2+\frac{\omega^2}{f'(r_h)^2}=0,
    \qquad
    \sigma_{\pm}=\pm\frac{i\omega}{f'(r_h)},
    \label{eq:indicial_horizon_sads}
\end{equation}
dengan turunan fungsi metrik pada horizon diberikan oleh
\begin{equation}
    f'(r_h)=\frac{2M}{r_h^2}+\frac{2r_h}{L^2}
    =\frac{1}{r_h}+\frac{3r_h}{L^2},
    \label{eq:fprime_sads}
\end{equation}
Sesuai dengan syarat batas gelombang masuk (\textit{ingoing}) pada horizon peristiwa, dipilih akar indikal negatif sehingga diperoleh
\begin{equation}
\sigma=-\frac{i\omega}{f'(r_h)}.
\end{equation}

Sementara itu, di batas tak hingga ($r\rightarrow\infty$), kondisi batas digunakan pada
ruang-waktu Schwarzschild Anti-de Sitter adalah kondisi batas Dirichlet, yaitu medan skalar
harus musnah pada batas konformal AdS ($\Psi\rightarrow0$ ketika $r\rightarrow\infty$)
\parencite{Horowitz2000,MossNorman2002}. Berbeda dengan ruang-waktu Schwarzschild yang
memiliki titik singular tak reguler di tak hingga, titik $r\rightarrow\infty$ pada SAdS
merupakan titik singular reguler dari persamaan radial, sehingga metode Frobenius standar
tetap berlaku. Akan tetapi, syarat $\Psi\rightarrow0$ saja belum menentukan laju peluruhan
solusi terhadap $r$, sehingga masih terdapat lebih dari satu kemungkinan perilaku asimtotik
yang konsisten dengan syarat tersebut. Untuk menentukan laju tersebut secara presisi,
perilaku asimtotik solusi ditinjau melalui persamaan indikal dengan mengasumsikan solusi
berbentuk $\Psi\sim r^\Delta$, sehingga diperoleh
\begin{equation}
    \Delta^2+\Delta-2=0, \qquad \Delta=1 \quad\text{atau}\quad \Delta=-2
    \label{eq:indicial_ads}
\end{equation}

Kedua akar tersebut merepresentasikan dua perilaku solusi di batas AdS, yaitu solusi yang
tumbuh sebanding dengan $r$ dan solusi yang meluruh sebanding dengan $r^{-2}$. Dari kedua
kemungkinan ini, hanya $\Delta=-2$ yang benar-benar memenuhi syarat $\Psi\rightarrow0$,
sedangkan $\Delta=1$ justru menyebabkan solusi divergen. Dengan demikian, kondisi batas
Dirichlet secara tegas memilih akar $\Delta=-2$.

Dengan menerapkan syarat batas Dirichlet tersebut, solusi Frobenius tidak memerlukan
\textit{prefaktor} tambahan di batas tak hingga dan cukup memfaktorkan perilaku singular di
sekitar horizon peristiwa. Hal ini karena variabel deret $z=1-r_h/r$ memetakan tak hingga
spasial ke $z=1$, sehingga perilaku peluruhan $\Delta=-2$ tidak perlu dituliskan sebagai
faktor pangkat $r^{-2}$ secara eksplisit, melainkan cukup dituntut melalui syarat bahwa
deret $\sum a_n z^n$ itu sendiri bernilai nol tepat di $z=1$. Bentuk solusi yang digunakan
adalah
\begin{equation}
    \Psi(r)=(r-r_h)^{-i\omega/f'(r_h)}\sum_{n=0}^{\infty} a_n z^n, \qquad z=1-\frac{r_h}{r}.
    \label{eq:ansatz_sads}
\end{equation}
dengan syarat Dirichlet dipenuhi melalui kondisi bahwa deret tersebut bernilai nol ketika
dievaluasi pada batas AdS, yaitu $z=1$ \parencite{Horowitz2000}.